% Ad Atticum 6.4 --- headnote
% ad-atticum-06-04, just after the Nones of June (~5 June) 50 BC, on the road to Tarsus

\textit{Cicero to Atticus, written on the march in
Cilicia just after the Nones of June (5 June) 50 BC ---
the manuscript dateline: \textit{Scr. in itinere paulo
post Non. Iun. a. 704 (50)}. A short, hurried note
written between halts as he comes down toward Tarsus
to close out his proconsular year. Atticus is by now
back at Rome; the letter is the first of the journey
sequence in which Cicero, no longer master of his time,
keeps reaching for his correspondent through the
disordered intervals of a march.}

\textit{Two anxieties run through the two sections.
The first is administrative and political: who will
hold Cilicia when Cicero lays it down. The senatorial
ruling that someone must be left in charge forces a
choice between an inadequate quaestor (Mescinius), a
Coelius from whom no news has come, and brother
Quintus --- the last preferred but freighted with the
costs of a separation and a war. The second is private
and confidential: Tullia's marriage (already settled
with Terentia in Atticus's absence and now requiring
his oversight), Cicero's own \textit{honor} --- the
supplicatio and the official despatch he needs Atticus
to have steered through the Senate --- and, in a
deliberately encrypted passage in Greek, what looks to
be a defalcation by Terentia's freedman in the matter
of property purchases. He cannot say more in writing;
the request is that Atticus, with his usual nose,
should follow the scent.}
